\section{Background \& Related Work}

\subsection{Provenance-Based IDS}
Provenance graphs represent system history by capturing relationships between system entities (processes, files, network sockets). Modern provenance IDS (e.g., OCR-APT \cite{ocrapt2025}, TraceCluster \cite{tracecluster2026}, StageFinder \cite{stagefinder2026}) construct subgraphs and apply Graph Neural Networks (GNNs) or sequence-based autoencoders to detect anomalies. However, all these models rely on offline, benign-only training datasets. If an attacker contaminates the training set, these models learn to classify the malicious behavior as normal.

Quantitatively, \textbf{OCR-APT} \cite{ocrapt2025} reports an APT detection AUC of 0.91 on DARPA Trace logs but requires 24--48 hours of clean pre-collection; \textbf{TraceCluster} \cite{tracecluster2026} achieves an average precision of 0.88 on multi-hop APT scenarios but uses a GNN trained on labeled subgraph clusters; \textbf{StageFinder} \cite{stagefinder2026} reports 94\% recall on staged APT campaigns but with a 12\% FPR under adversarial noise. In contrast, our hybrid framework leverages SCM causal invariants to build robust baseline models directly on raw streams, using a small, partially labeled validation partition to train and fuse ML components rather than assuming fully clean offline training data.

\subsection{Unsupervised and Reduced-Label Provenance IDS}
To relax the labeling burden, several systems attempt unsupervised or self-supervised threat detection:
\begin{itemize}
\item \textbf{UNICORN} (2020) \cite{han2020unicorn} summarizes streaming provenance subgraphs as graph histograms and fits an unsupervised clustering baseline. However, it requires a clean, attack-free training period to define normal system execution clusters.
\item \textbf{WATSON} (2020) \cite{gao2020watson} aggregates audit events and models behaviors using contextual semantics. While reducing manual modeling, it relies on static, clean offline profiles of normal execution and has no mechanism to adapt to concept-drift.
\item \textbf{ShadeWatcher} (2022) \cite{yang2022shadewatcher} frames threat analysis as a recommendation task to flag anomalous edges. It requires a clean audit log baseline and leverages user feedback (labels) to guide its recommendation engine.
\item \textbf{Kairos} (2024) \cite{cheng2024kairos} utilizes self-supervised learning over provenance graphs to minimize label dependency, yet it still depends on an uncontaminated pre-collection phase for initial parameter fitting.
\end{itemize}
Unlike these approaches, our semi-supervised model operates on raw, uncurated streams. By using a small validation slice with sparse labels to calibrate the meta-learner and the Behavioral Pattern Reasoner, it bridges the gap between purely unsupervised models (which lack alignment with actual attack patterns) and supervised models (which require massive labeled data).

\subsection{Causal State Inference}
Causal reasoning is increasingly applied to security to capture system mechanics. \textbf{Causal-IDS} (2026) \cite{causalids2026} models network flow log variables using static SCMs to identify intrusions as violations of causal mechanisms. While sharing our focus on causal violations, Causal-IDS differs fundamentally from our work:
\begin{itemize}
\item \textbf{Network vs. Provenance}: Causal-IDS operates on network flow statistics (e.g., packet counts, byte rates), whereas our system focuses on fine-grained provenance-graph edges to detect APT behaviors.
\item \textbf{Benign Data Reliance}: Causal-IDS requires clean training logs to fit its SCM, whereas our framework learns causal structures online from unlabeled, potentially contaminated streams.
\item \textbf{Behavioral State Reasoning}: Causal-IDS relies purely on causal p-values, which are sensitive to estimation error. Our framework projects SCM outputs into a multi-dimensional behavioral state space, applying local, temporal, and behavioral reasoning operators to achieve far higher classification performance.
\item \textbf{Thresholding}: Causal-IDS uses static, empirical thresholds, whereas our framework provides distribution-free false-alarm bounds via online conformal prediction.
\end{itemize}

Other systems, such as \textbf{CausalGraph} \cite{causalgraph2026}, rely on LLMs to perform causal reasoning over provenance subgraphs, which is computationally expensive and slow for real-time high-velocity logs.

\subsection{Comparative Summary}
Table~\ref{tab:related_comparison} situates our framework against representative provenance-based IDS methods across four operational dimensions. The ``Label-Free Train'' column indicates whether the model can learn a normal baseline from fully unlabeled, potentially contaminated logs; ``Drift-Adaptive'' indicates whether the model adjusts to natural concept drift at runtime; ``Causal'' indicates whether causal structure is explicitly modeled; and ``FPR Guarantee'' indicates whether a formal bound on the false alarm rate is provided.

\begin{table*}[t]
\centering
\small
\caption{Comparison of provenance IDS methods on key deployment properties. ``Clean Baseline'' = requires a clean, attack-free pre-collection period. Performance numbers are from respective publications and are \emph{not} directly comparable across rows (different datasets/attacker models).}
\begin{tabularx}{\textwidth}{Xccccl}
\toprule
Method & \makecell{Clean\\Baseline} & \makecell{Drift\\Adaptive} & Causal & \makecell{FPR\\Guarantee} & Best Reported \\
\midrule
OCR-APT \cite{ocrapt2025}            & Yes (48h)   & No  & No  & No  & AUC 0.91 (DARPA Trace) \\
TraceCluster \cite{tracecluster2026} & Yes (GNN)   & No  & No  & No  & AP 0.88 (multi-hop APT) \\
StageFinder \cite{stagefinder2026}   & Yes         & No  & No  & No  & Recall 94\% / FPR 12\% \\
Kairos \cite{cheng2024kairos}        & Yes         & No  & No  & No  & F1 0.89 (DARPA Trace) \\
UNICORN \cite{han2020unicorn}        & Yes         & No  & No  & No  & --- \\
Causal-IDS \cite{causalids2026}      & Yes         & No  & Yes & No  & --- \\
\textbf{Ours (CBSR)}                 & \textbf{No} & \textbf{Yes} & \textbf{Yes} & \textbf{Yes} & \textbf{AUC 0.9628 (OpTC)} \\
\bottomrule
\end{tabularx}
\label{tab:related_comparison}
\end{table*}

\noindent\textit{Scope note:} The methods above are evaluated on different datasets; direct AUC comparison is not meaningful across rows. The table highlights \emph{capability gaps} rather than performance rankings. Our framework is the only method that simultaneously operates without a clean baseline, adapts to drift, models causal structure, and provides a formal FPR bound.
